% Zone „Das kennst du schon“ – aus bank/lineare-gleichungssysteme/zone.jsonl (zusammenbau v0.1)
\einheitenkopf[zone][Kennst du schon]{Das kennst du schon}
\setcounter{aufgabe}{0}
%% TODO Grundvorstellungs-Aufgabe als erste Hauptnummer der Zone (2.8) fehlt in der Bank

%% TODO Titel: Ich-kann-Satz fehlt in der Bank (Platzhalter: Kettenname) – Nr. 1
%% TODO Anweisung: ein Satz über der ersten Teilaufgabe fehlt in der Bank – Nr. 1
\begin{aufgabe}{Lösung durch Einsetzen prüfen, Ergebnis mit (wA)/(fA)}
%% TODO Darstellung für schwach fehlt in der Bank (lineare-gleichungssysteme-zone-f1-v1; Abschnitt „Für schwache Schüler“)
\swz{$x = 4$ – Lösung von $x + 5 = 9$? Setze ein.}{}
%% TODO Darstellung für schwach fehlt in der Bank (lineare-gleichungssysteme-zone-f1-v3; Abschnitt „Für schwache Schüler“)
\swz{$x = -2$ – Lösung von $3x + 14 = 8$? Setze ein.}{}
\end{aufgabe}

%% TODO Titel: Ich-kann-Satz fehlt in der Bank (Platzhalter: Kettenname) – Nr. 2
%% TODO Anweisung: ein Satz über der ersten Teilaufgabe fehlt in der Bank – Nr. 2
\begin{aufgabe}{Gleichung nach einer Variablen umstellen (Formel umstellen: ax + by = c nach y)}
%% TODO Darstellung für schwach fehlt in der Bank (lineare-gleichungssysteme-zone-f2-v1; Abschnitt „Für schwache Schüler“)
\swz{$\text{$x + y = 15$ – nach $y$ umstellen.}$}{}
%% TODO Darstellung für schwach fehlt in der Bank (lineare-gleichungssysteme-zone-f2-v3; Abschnitt „Für schwache Schüler“)
\swz{$\text{$4x + 2y = 14$ – nach $y$ umstellen.}$}{}
\end{aufgabe}

%% TODO Titel: Ich-kann-Satz fehlt in der Bank (Platzhalter: Kettenname) – Nr. 3
%% TODO Anweisung: ein Satz über der ersten Teilaufgabe fehlt in der Bank – Nr. 3
\begin{aufgabe}{Gerade aus der Gleichung zeichnen (n und Steigungsdreieck), Punkt ablesen, Lage zweier Geraden (parallel bei gleicher Steigung)}
\swa{Gerade $y = 2x + 1$ – zeichne sie mit Achsenabschnitt und Steigungsdreieck.}{\begin{ksys}[xmin=-5,xmax=5,ymin=-5,ymax=5] \end{ksys}}
\swa{Gerade $y = -\frac{1}{2}x + 3$ – zeichne sie mit Achsenabschnitt und Steigungsdreieck.}{\begin{ksys}[xmin=-5,xmax=5,ymin=-5,ymax=5] \end{ksys}}
\end{aufgabe}

%% TODO Titel: Ich-kann-Satz fehlt in der Bank (Platzhalter: Kettenname) – Nr. 4
%% TODO Anweisung: ein Satz über der ersten Teilaufgabe fehlt in der Bank – Nr. 4
\begin{aufgabe}{Lineare Gleichung mit Klammern, Dezimalzahlen und x auf beiden Seiten lösen, Schreibform mit Strich}
%% TODO Darstellung für schwach fehlt in der Bank (lineare-gleichungssysteme-zone-f4-v1; Abschnitt „Für schwache Schüler“)
\swz{$2x + 5 = 17$}{}
%% TODO Darstellung für schwach fehlt in der Bank (lineare-gleichungssysteme-zone-f4-v3; Abschnitt „Für schwache Schüler“)
\swz{$0{,}5x + 2 = 4{,}5$}{}
\end{aufgabe}

%% TODO Titel: Ich-kann-Satz fehlt in der Bank (Platzhalter: Kettenname) – Nr. 5
%% TODO Anweisung: ein Satz über der ersten Teilaufgabe fehlt in der Bank – Nr. 5
\begin{aufgabe}{Klammer ausmultiplizieren (Zahl mal Klammer, Minus vor der Klammer), Terme mit zwei Variablen zusammenfassen}
%% TODO Darstellung für schwach fehlt in der Bank (lineare-gleichungssysteme-zone-f5-v1; Abschnitt „Für schwache Schüler“)
\swz{$3(x + 4)$ – ausmultiplizieren.}{}
%% TODO Darstellung für schwach fehlt in der Bank (lineare-gleichungssysteme-zone-f5-v3; Abschnitt „Für schwache Schüler“)
\swz{$4(2x - 3y)$ – ausmultiplizieren.}{}
\end{aufgabe}

%% TODO Titel: Ich-kann-Satz fehlt in der Bank (Platzhalter: Kettenname) – Nr. 6
%% TODO Anweisung: ein Satz über der ersten Teilaufgabe fehlt in der Bank – Nr. 6
\begin{aufgabe}{Rechnen mit Dezimalzahlen (Geld) und negativen Zahlen, Division mit dem Taschenrechner}
%% TODO Darstellung für schwach fehlt in der Bank (lineare-gleichungssysteme-zone-f6-v1; Abschnitt „Für schwache Schüler“)
\swz{$3 \cdot 2{,}40\,€$ – wie viel?}{}
%% TODO Darstellung für schwach fehlt in der Bank (lineare-gleichungssysteme-zone-f6-v3; Abschnitt „Für schwache Schüler“)
\swz{$4 \cdot 1{,}35\,€ + 2 \cdot 0{,}80\,€$ – wie viel zusammen?}{}
\end{aufgabe}

%% TODO Titel: Ich-kann-Satz fehlt in der Bank (Platzhalter: Kettenname) – Nr. 7
%% TODO Anweisung: ein Satz über der ersten Teilaufgabe fehlt in der Bank – Nr. 7
\begin{aufgabe}{Gleichung mit einer Variablen aus einem Sachverhalt aufstellen (Preis · Anzahl, „zusammen“)}
%% TODO Darstellung für schwach fehlt in der Bank (lineare-gleichungssysteme-zone-f7-v1; Abschnitt „Für schwache Schüler“)
\swz{$\text{Drei Hefte kosten zusammen $4{,}50\,€$. Preis eines Hefts? Stelle eine Gleichung für den Preis $x$ eines Hefts auf und löse sie.}$}{}
%% TODO Darstellung für schwach fehlt in der Bank (lineare-gleichungssysteme-zone-f7-v3; Abschnitt „Für schwache Schüler“)
\swz{$\text{Ein Kinoticket kostet $x$ Euro. Vier Tickets und eine Popcorntüte für $3{,}60\,€$ kosten zusammen $38{,}00\,€$. Preis eines Tickets? Stelle eine Gleichung auf und bestimme $x$.}$}{}
\end{aufgabe}

%% TODO Titel: Ich-kann-Satz fehlt in der Bank (Platzhalter: Kettenname) – Nr. 8
%% TODO Anweisung: ein Satz über der ersten Teilaufgabe fehlt in der Bank – Nr. 8
\begin{aufgabe}{Sek II: Terme mit drei Variablen zusammenfassen, Gleichungen vervielfachen und addieren, nach einer Variablen umstellen}
%% TODO Darstellung für schwach fehlt in der Bank (lineare-gleichungssysteme-zone-f8-v1; Abschnitt „Für schwache Schüler“)
\swz{$(x + 2y - z) + (3x - 2y + 4z)$ – zusammenfassen.}{}
%% TODO Darstellung für schwach fehlt in der Bank (lineare-gleichungssysteme-zone-f8-v3; Abschnitt „Für schwache Schüler“)
\swz{I: $x + 3y - 2z = 4$, II: $-x + y + 5z = 7$ – addiere I und II.}{}
\end{aufgabe}

%% TODO Titel: Ich-kann-Satz fehlt in der Bank (Platzhalter: Kettenname) – Nr. 9
%% TODO Anweisung: ein Satz über der ersten Teilaufgabe fehlt in der Bank – Nr. 9
\begin{aufgabe}{Sek II: lineare Ungleichungen in einer Variablen lösen und Lösungen sieben (Vorzeichenbedingung, Ganzzahligkeit, größter Wert)}
%% TODO Darstellung für schwach fehlt in der Bank (lineare-gleichungssysteme-zone-f9-v1; Abschnitt „Für schwache Schüler“)
\swz{$x + 3 < 9$ – welche Zahlen $x$?}{}
%% TODO Darstellung für schwach fehlt in der Bank (lineare-gleichungssysteme-zone-f9-v3; Abschnitt „Für schwache Schüler“)
\swz{$3t - 4 \le 14$ – welche Zahlen $t$? Gib die größte ganze Zahl an.}{}
\end{aufgabe}

%% TODO Titel: Ich-kann-Satz fehlt in der Bank (Platzhalter: Kettenname) – Nr. 10
%% TODO Anweisung: ein Satz über der ersten Teilaufgabe fehlt in der Bank – Nr. 10
\begin{aufgabe}{Sek II: Tripel als Punkte und Vektoren lesen und die Lösungsmenge als Menge schreiben}
%% TODO Darstellung für schwach fehlt in der Bank (lineare-gleichungssysteme-zone-f10-v1; Abschnitt „Für schwache Schüler“)
\swz{Punkt $P(2 \mid -1 \mid 4)$ – welche Zahl ist die $y$-Koordinate?}{}
%% TODO Darstellung für schwach fehlt in der Bank (lineare-gleichungssysteme-zone-f10-v3; Abschnitt „Für schwache Schüler“)
\swz{Lösungsmenge $\{(1 + t;\ 2t;\ t) \mid t \in \mathbb{R}\}$ – welches Tripel gehört zu $t = 3$?}{}
\end{aufgabe}

%% TODO Titel: Ich-kann-Satz fehlt in der Bank (Platzhalter: Kettenname) – Nr. 11
%% TODO Anweisung: ein Satz über der ersten Teilaufgabe fehlt in der Bank – Nr. 11
\begin{aufgabe}{Klammer ausmultiplizieren (Zahl mal Klammer, Minus vor der Klammer), Terme mit zwei Variablen zusammenfassen – Fehler finden}
\swz[3]{Tom rechnet: \rechnung{14 - (y + 5) &= 14 - y + 5 \\ &= 19 - y} Finde den Fehler.}{}
\end{aufgabe}

%% TODO Titel: Ich-kann-Satz fehlt in der Bank (Platzhalter: Kettenname) – Nr. 12
%% TODO Anweisung: ein Satz über der ersten Teilaufgabe fehlt in der Bank – Nr. 12
\begin{aufgabe}{Klammer ausmultiplizieren (Zahl mal Klammer, Minus vor der Klammer), Terme mit zwei Variablen zusammenfassen – selbst rechnen}
%% TODO Darstellung für schwach fehlt in der Bank (lineare-gleichungssysteme-zone-f5-v6; Abschnitt „Für schwache Schüler“)
\swz{$16 - (x + 9)$ – Klammer auflösen und zusammenfassen.}{}
\end{aufgabe}
